\documentclass[12pt,a4paper]{article}
\usepackage[utf8]{inputenc}
\usepackage[margin=1in]{geometry}
\usepackage{amsmath}
\usepackage{amssymb}
\usepackage{graphicx}
\usepackage{xcolor}
\usepackage{hyperref}
\usepackage{fancyhdr}
\usepackage{array}
\usepackage{booktabs}

\pagestyle{fancy}
\lhead{Estratégia Adaptativa Multi-Fator}
\rhead{Etapa 3: Backtest}
\cfoot{\thepage}

\title{\textbf{METODOLOGIA E SETUP DE BACKTEST}\\[0.5cm]
\textit{Adaptive Multi-Factor Regime-Switching Strategy}}
\author{Quantitative Research}
\date{\today}

\begin{document}

\maketitle

\section{Dados Utilizados}

\subsection{Período de Análise}

\begin{itemize}
    \item \textbf{Data de Início}: 3 anos passados
    \item \textbf{Data de Fim}: Hoje
    \item \textbf{Frequência}: Diária (252 dias de trading por ano)
    \item \textbf{Total de observações}: 756 dias de trading
\end{itemize}

\subsection{Universo de Ativos}

\begin{itemize}
    \item \textbf{Número de ativos}: 50 ações
    \item \textbf{Mercado}: Mercado de ações líquido
    \item \textbf{Critério de seleção}: Ativos com histórico completo de 3 anos
\end{itemize}

\subsection{Dados Fundamentalistas}

Coletamos os seguintes dados de cada ativo:
\begin{itemize}
    \item Book Value (valor contábil)
    \item ROE (Return on Equity)
    \item Profit Margin (margem de lucro)
    \item Revenue Growth (crescimento de receita)
    \item Total Assets
    \item EBIT
\end{itemize}

Frequência: dados fundamentalistas trimestrais (forward-filled para daily).

\section{Configuração do Backtest}

\subsection{Parâmetros Iniciais}

\begin{table}[h]
\centering
\begin{tabular}{|l|c|}
\hline
\textbf{Parâmetro} & \textbf{Valor} \\
\hline
Capital Inicial & \$1.000.000 \\
Rebalanceamento & Mensal (a cada 21 dias de trading) \\
Lookback Window & 252 dias (1 ano) \\
Limite de alavancagem & Dinâmico por regime \\
Limite diário de loss & -10\% por dia \\
Custo de transação & 10 bps (0.001\%) por rebalanceamento \\
Slippage & 5 bps \\
\hline
\end{tabular}
\end{table}

\subsection{Warm-up Period}

Primeira ano de dados (252 dias) utilizado apenas para:
\begin{itemize}
    \item Estimar matriz de covariância
    \item Treinar HMM inicial
    \item Calibrar modelos ML
\end{itemize}

Backtest inicia em $t = 252$.

\section{Metodologia de Análise}

\subsection{Daily Rebalancing Logic}

A cada dia $t$:

\begin{enumerate}
    \item \textbf{Factor Computation}: Calcular 7 fatores usando dados históricos até $t-1$
    
    \item \textbf{Regime Detection}: Aplicar HMM em últimas 252 observações para identificar $s_t$
    
    \item \textbf{Rebalancing Check}: Se $(t - 252) \mod 21 = 0$:
    \begin{enumerate}
        \item Treinar GBM e Ridge em dados históricos $[t-252, t-20]$
        \item Prever retornos esperados $\hat{\mu}_t$
        \item Estimar matriz de covariância $\hat{\Sigma}_t$ em última janela de 252 dias
        \item Otimizar portfolio para pesos $w_t^*$
        \item Executar trades de rebalanceamento
    \end{enumerate}
    
    \item \textbf{Daily Returns}: Computar retorno diário realizado
    \begin{align}
        r_{portfolio,t} = w_t^T \cdot r_{t}
    \end{align}
    onde $r_t$ é vetor de retornos dos ativos no dia $t$.
    
    \item \textbf{Atualizar Value}: $V_t = V_{t-1} \cdot (1 + r_{portfolio,t})$
\end{enumerate}

\subsection{Gestão de Risco}

\subsubsection{Limite de Loss Diário}

Se retorno diário < -10\%, aplicamos stop-loss. Portfolio retorna para 70\% de alocação (posição defensiva).

\[
r_{portfolio,t} = \begin{cases}
r_{portfolio,t} & \text{se } r_{portfolio,t} \geq -10\% \\
-10\% & \text{se } r_{portfolio,t} < -10\%
\end{cases}
\]

\subsubsection{Drawdown Management}

Quando drawdown acumulado > 20\%, reduzir leverage para 0.5.

\[
\text{drawdown}_t = \frac{V_t - V_{\max}}{V_{\max}}
\]

\subsubsection{Volatility Targeting}

Pesos ajustados diariamente para manter volatilidade alvo de 8\% a.a.:

\[
w_{\text{final}, t} = w_{\text{optimized}, t} \cdot \frac{0.08}{\sqrt{252} \cdot \hat{\sigma}_t}
\]

onde $\hat{\sigma}_t$ é volatilidade estimada em última janela de 20 dias.

\section{Benchmarks de Comparação}

Estratégia é comparada contra 3 benchmarks:

\subsection{Benchmark 1: Buy-and-Hold S\&P 500}

Portfolio 100\% investido em índice de mercado. Retorno e risco baseline.

\subsection{Benchmark 2: 60/40 Balanced Portfolio}

Alocação conservadora: 60\% ações, 40\% títulos (fixed income com taxa de 2\% a.a.).

\subsection{Benchmark 3: Equal-Weight Portfolio}

Peso igual em todos 50 ativos, rebalanceado mensalmente. Controle simples.

\section{Métricas de Performance}

\subsection{Retorno}

\[
\text{Total Return} = \frac{V_T - V_0}{V_0}
\]

\[
\text{Annual Return} = \left(\frac{V_T}{V_0}\right)^{252/T} - 1
\]

\subsection{Volatilidade}

\[
\sigma = \sqrt{\frac{1}{T-1} \sum_{t=1}^{T} (r_t - \bar{r})^2} \cdot \sqrt{252}
\]

Volatilidade anualizada dos retornos diários.

\subsection{Sharpe Ratio}

\[
\text{Sharpe} = \frac{\text{Annual Return} - R_f}{\sigma} 
\]

com $R_f = 2\%$ (risk-free rate).

\subsection{Sortino Ratio}

Penaliza apenas downside volatility:

\[
\text{Sortino} = \frac{\text{Annual Return} - R_f}{\sigma_{\text{downside}}}
\]

onde $\sigma_{\text{downside}} = \sqrt{\frac{1}{T}\sum_{r_t < 0} r_t^2}$.

\subsection{Calmar Ratio}

\[
\text{Calmar} = \frac{\text{Annual Return}}{|\text{Max Drawdown}|}
\]

Mostra retorno por unidade de risco de queda máxima.

\subsection{Maximum Drawdown}

\[
\text{Max DD} = \min_t \frac{V_t - V_{\max_{s \leq t}}}{V_{\max_{s \leq t}}}
\]

Pior queda de pico a vale durante período.

\subsection{Value-at-Risk (VaR) 95\%}

\[
\text{VaR}_{95\%} = \text{Percentil}_{5\%}(r_1, \ldots, r_T)
\]

Worst 5\% dos dias de trading (pior loss esperado).

\subsection{Conditional VaR (CVaR) 95\%}

\[
\text{CVaR}_{95\%} = E[r_t \mid r_t \leq \text{VaR}_{95\%}]
\]

Média dos piores 5\% dos dias. Mais relevante que VaR para tail risk.

\subsection{Win Rate}

\[
\text{Win Rate} = \frac{\#\{\text{dias com } r_t > 0\}}{\text{Total de dias}}
\]

Percentual de dias com retorno positivo.

\subsection{Profit Factor}

\[
\text{Profit Factor} = \frac{\sum_{r_t > 0} r_t}{|\sum_{r_t < 0} r_t|}
\]

Razão de ganhos acumulados vs. perdas acumuladas.

\section{Tratamento de Custos de Transação}

\subsection{Custo de Rebalanceamento}

A cada rebalanceamento, aplicamos custo de 10 bps:

\[
C_t = 10 \text{ bps} \times \sum_i |w_{i,t} - w_{i,t-1}|
\]

Este custo é deduzido do retorno diário.

\subsection{Market Slippage}

Adicionamos 5 bps de slippage em cada ordem de compra/venda, modelando impacto de mercado.

\section{Validação de Viés}

\subsection{Look-Ahead Bias}

Garantimos que nenhuma informação do futuro é usada:
\begin{itemize}
    \item Fatores computados com dados até $t-1$ apenas
    \item ML treinado em janelas pasadas ($[t-252, t-20]$)
    \item Fundamentals forward-filled (não usando dados futuros)
\end{itemize}

\subsection{Survivorship Bias}

Backtest inclui ativos que podem deixar o universo, simulando real-world conditions.

\subsection{Data Snooping}

Modelo foi especificado teoricamente antes de backtest. Parâmetros não foram otimizados via grid search (evitar overfitting).

\section{Cenários de Teste}

\subsection{Base Case}

Backtest padrão com parâmetros nominais (como descrito acima).

\subsection{Stress Test 1: Aumento de Volatilidade}

Simular ambiente de alta volatilidade (tipo COVID 2020):
\begin{itemize}
    \item Multiplicar volatilidades históricas por 2x
    \item Correlações aumentam para 0.7-0.9 entre ativos
    \item Testar resiliência de hedge do modelo
\end{itemize}

\subsection{Stress Test 2: Crash de Mercado}

Simular crash tipo 2008:
\begin{itemize}
    \item Retorno de -20\% no mercado em 1 dia
    \item Testar se HMM detecta mudança de regime rápido
    \item Verificar se portfolio se defende adequadamente
\end{itemize}

\end{document}
